\section{Revisión de medio ambiente y salud}

\subsection{Estudio de impactos en la salud por contaminantes}

El análisis de la calidad del aire es importante para la ciudadanía, porque la
contaminación atmosférica representa un riesgo para la salud y se relaciona
con enfermedades como cardiopatías, accidentes cerebrovasculares, neumonía,
cáncer de pulmón e infecciones respiratorias. Este estudio analiza diversos
contaminantes \parencite{oms2021}:

\begin{itemize}
	\item partículas suspendidas \ce{PM2.5} y \ce{PM10};
	\item dióxido de nitrógeno (\ce{NO2});
	\item ozono (\ce{O3}); y
	\item dióxido de azufre (\ce{SO2}).
\end{itemize}

También considera sus concentraciones, la exposición de la población y los
datos de salud. La calidad del aire se mide mediante estaciones de monitoreo
que registran concentraciones horarias, a partir de las cuales se calculan
métricas como promedios de 24 horas o anuales, dependiendo del contaminante.
El estudio utiliza normas oficiales mexicanas, como la NOM-025-SSA1-2021 para
\ce{PM2.5}, y las compara con las recomendaciones de la Organización Mundial
de la Salud (OMS).

Una problemática identificada es que algunos límites mexicanos son menos
estrictos que los recomendados por la OMS. Por ejemplo, para \ce{PM2.5}, el
límite mexicano es de $10\,\mu\mathrm{g}/\mathrm{m}^{3}$, mientras que la OMS
recomienda $5\,\mu\mathrm{g}/\mathrm{m}^{3}$. Además, la medición enfrenta
dificultades porque no todas las estaciones miden todos los contaminantes ni
cuentan con suficientes datos válidos. Por ello, se requieren criterios de
suficiencia, como contar con al menos el 75\% de datos válidos
\parencite{nom025ssa1,oms2021}.

\subsection{Contaminación del aire doméstico: riesgos relativos}

Este documento informa sobre los efectos de la contaminación del aire en la
salud. Presenta riesgos importantes y su relación con diferentes enfermedades,
entre ellas infecciones respiratorias de vías aéreas inferiores, enfermedad
pulmonar obstructiva crónica, cáncer de pulmón, cardiopatía isquémica y
accidente cerebrovascular. También considera diferencias según la edad y el
sexo, y utiliza intervalos de confianza del 95\% para expresar la incertidumbre
de los resultados.

Por ejemplo, el riesgo relativo es particularmente elevado para el cáncer de
pulmón en mujeres. En conclusión, el documento muestra que analizar la
contaminación del aire es importante para conocer y cuantificar sus posibles
efectos sobre la salud de la población.

\subsection{Normativa de aire de Nuevo León}

Esta fuente permite conocer las normas oficiales y la interpretación de las
mediciones de calidad del aire. Las normas establecen límites de concentración
para diferentes contaminantes y consideran tanto la cantidad del contaminante
como el tiempo de exposición, ya que los límites pueden establecerse para
periodos de una hora, ocho horas, 24 horas o promedios anuales.

Por ejemplo, la Norma Oficial Mexicana NOM-021-SSA1-2021 evalúa la calidad del
aire ambiente con respecto al monóxido de carbono (\ce{CO}). La página reúne
12 normas oficiales mexicanas
\parencite{nlAirenlgobmxNormatividad,NOMCalidadAireAmbiente}.

Las mediciones permiten comparar la concentración registrada con los valores
establecidos en las normas para determinar si se cumplen los límites
correspondientes. La fuente confirma que los límites se establecen considerando
los posibles efectos sobre la salud, pero no desarrolla una clasificación
internacional de contaminantes.
